\chapter{État de l'art et choix technologiques}

\section*{Introduction}
\addcontentsline{toc}{section}{Introduction}

Ce chapitre rassemble les notions sur lesquelles notre solution s'appuie. Nous y
revenons sur le DevOps, la chaîne CI/CD et leur extension DevSecOps, sur les grands
modèles de langage et l'inférence locale, sur l'orchestration des workflows IA, ainsi
que sur les outils d'analyse de sécurité statique et la stack applicative. Une étude
comparative des technologies disponibles pour chaque composant permet ensuite de
justifier les choix retenus. Le chapitre se clôt sur la définition des besoins
fonctionnels et non fonctionnels du projet.

% ─────────────────────────────────────────────────────────────────────────────
\section{DevOps, CI/CD et DevSecOps}

    \subsection{Le DevOps}

        Le DevOps regroupe un ensemble de pratiques, d'outils et de choix culturels qui
        cherchent à raccourcir le cycle de livraison logicielle tout en gardant le produit
        stable. Il s'organise autour de trois axes : la collaboration au quotidien entre
        développement et exploitation, là où le modèle classique séparait les deux ;
        l'automatisation des tâches répétitives de la chaîne (build, tests, déploiement)
        pour limiter l'erreur humaine et libérer du temps ; et le suivi continu en
        production, qui permet de détecter rapidement les régressions et d'alimenter les
        itérations suivantes.

    \subsection{L'approche CI/CD}

        \subsubsection{Définition}

        L'approche CI/CD automatise les étapes de build, de test et de déploiement, ce
        qui raccourcit le délai entre le développement et la mise en production tout en
        réduisant les risques d'intégration \cite{DevSecOps}. Elle se décompose en deux
        phases enchaînées : l'intégration continue côté code et le déploiement continu
        côté production.

        \subsubsection{Intégration continue (CI)}

        L'intégration continue implique de modifier régulièrement le code de l'application,
        de le tester automatiquement, puis de le fusionner dans un référentiel commun. Dans
        notre projet, chaque ouverture ou mise à jour d'une \textit{Pull Request} sur GitHub
        déclenche automatiquement le pipeline de sécurité de l'agent via le mécanisme de
        webhooks.

        \subsubsection{Déploiement continu (CD)}

        Le déploiement continu automatise l'étape suivante du pipeline après l'intégration
        continue, en permettant la mise en production des modifications validées sans
        intervention manuelle. Dans notre contexte, le verdict de l'agent (APPROVE /
        REQUEST\_CHANGES / BLOCK) sert de gate de qualité et de sécurité qui conditionne
        la fusion automatique ou bloquée de la \textit{Pull Request}.

    \subsection{Du DevOps au DevSecOps}

        Le DevOps tel qu'il vient d'être décrit traite la sécurité comme une étape
        séparée, généralement repoussée à la fin du cycle. Le DevSecOps, contraction
        de Development, Security et Operations \cite{DevSecOps}, corrige cette omission
        en intégrant la sécurité à chaque étape du cycle plutôt qu'en aval
        \cite{NIST_DevSecOps}.
        L'objectif est de détecter les vulnérabilités au plus tôt, là où elles coûtent le
        moins cher à corriger ; un bug de sécurité repéré sur une \textit{Pull Request}
        avant fusion n'a presque aucun coût comparé à la même faille découverte en
        production.

        \begin{figure}[H]
            \centering
            \frame{\includegraphics[width=0.82\columnwidth]{devsecops_lifecycle}}
            \caption{Cycle de vie DevSecOps — intégration de la sécurité à chaque étape. \\
                     \footnotesize\textit{Source : Rezadlt, Wikimedia Commons, CC BY-SA 4.0.}}
            \label{fig:devsecops}
        \end{figure}

        Comme le montre la figure \ref{fig:devsecops}, le DevSecOps s'organise autour de
        quelques principes structurants. Le \textit{Security as Code} consiste à exprimer
        les politiques de sécurité sous forme de code versionné et déployé automatiquement,
        au même titre qu'une configuration applicative. Le \textit{Shift Left} pousse la
        détection des vulnérabilités le plus en amont possible, idéalement sur la
        \textit{Pull Request} et avant toute fusion. Le \textit{feedback continu} complète
        l'ensemble en remontant les résultats des analyses aux développeurs concernés, sur
        leur outil de travail, dès qu'ils sont disponibles.

    \subsection{L'OWASP Top 10}

        L'OWASP (Open Web Application Security Project) publie périodiquement la liste des dix
        risques de sécurité les plus critiques pour les applications web \cite{OWASP2025}. La
        version 2025 sert de référence internationale et de grille d'analyse pour le moteur LLM
        de notre solution. Parmi les évolutions notables par rapport à la version 2021 figurent
        l'apparition de la catégorie A03:2025 --- \textit{Software Supply Chain Failures}, qui
        élargit l'ancien périmètre « composants vulnérables » à l'ensemble de la chaîne
        d'approvisionnement logicielle (dépendances, registres, pipelines CI/CD), et la nouvelle
        catégorie A10:2025 --- \textit{Mishandling of Exceptional Conditions} qui remplace SSRF.
        La catégorie A03:2025 recouvre directement le domaine couvert par OSV-Scanner et Trivy
        dans notre pipeline. Chaque finding produit par l'agent est catégorisé selon cette
        taxonomie avec une référence explicite (par exemple A01:2025 --- Broken Access Control
        ou A05:2025 --- Injection).

    \subsection{Webhooks GitHub et signature HMAC-SHA256}

        Le mécanisme concret qui relie ce cycle DevSecOps à notre agent est celui des
        \textit{webhooks} GitHub \cite{GitHubWebhooks}. À chaque ouverture ou mise à jour
        d'une \textit{Pull Request}, GitHub émet une requête HTTP POST vers l'endpoint
        public de l'agent, accompagnée d'une charge utile JSON décrivant l'événement
        (dépôt, branche, auteur, identifiant de la PR). Cet appel sortant est ce qui
        permet au pipeline de sécurité de se déclencher sans aucune action manuelle.

        Pour garantir que seul GitHub puisse déclencher le pipeline, chaque requête est
        signée par GitHub avec un secret partagé via l'algorithme HMAC-SHA256 \cite{HMAC}.
        L'agent recalcule la signature à partir de la charge utile reçue et la compare en
        temps constant à celle envoyée dans l'en-tête \texttt{X-Hub-Signature-256} : toute
        requête dont la signature ne correspond pas est rejetée en 401 avant la moindre
        opération. Ce double mécanisme associe un déclenchement automatique côté
        événement à une validation cryptographique côté réception. Il rejoint ainsi
        à la fois l'idée d'automatisation du DevOps et celle de \textit{Security as
        Code} du DevSecOps.

% ─────────────────────────────────────────────────────────────────────────────
\section{Les grands modèles de langage (LLM)}

    \subsection{Fonctionnement général}

        Les grands modèles de langage (Large Language Models, LLM) sont des réseaux de
        neurones de type Transformeur entraînés sur d'immenses corpus de texte et de code
        \cite{QwenCoder}.
        Ils sont capables de comprendre le contexte d'un fichier source et d'identifier
        des patterns de vulnérabilité connus, tels que la concaténation directe de
        chaînes dans une requête SQL ou l'utilisation de fonctions dangereuses de
        désérialisation. Ils savent également proposer des corrections dans le langage
        du développeur. Le paramètre \texttt{num\_ctx} (fenêtre de contexte) est
        critique dans notre cas. Un contexte insuffisant oblige le modèle à tronquer
        l'entrée ou à compléter depuis sa mémoire d'entraînement, ce qui produit des
        hallucinations. Nous avons calibré ce paramètre à
        12 288 tokens pour l'analyse de \textit{Pull Request} et à 6 144 tokens pour
        l'assistant de chat.

    \subsection{Inférence locale et API cloud}

        Deux paradigmes coexistent pour exécuter un LLM. L'\textbf{inférence locale}
        consiste à héberger le modèle sur l'infrastructure de l'organisation : le code
        soumis au modèle ne quitte jamais le périmètre interne, mais les performances
        dépendent du matériel disponible. L'\textbf{API cloud} consiste à invoquer un
        modèle exposé par un fournisseur tiers (OpenAI, Anthropic, etc.) : la puissance
        de calcul est déléguée, mais chaque requête transmet le code à des serveurs
        externes. Pour un établissement bancaire soumis aux exigences de confidentialité
        de la Banque Centrale de Tunisie, ce choix conditionne directement la faisabilité
        de la solution.

    \subsection{Tool calling}

        Le \textit{tool calling} (ou \textit{function calling}) désigne la capacité d'un
        LLM à émettre, en sortie, une demande structurée d'invocation d'une fonction
        externe au lieu d'une simple réponse textuelle. Le modèle reçoit en entrée la
        description d'un ensemble d'outils (nom, paramètres attendus, format de retour)
        et décide, en fonction de la requête de l'utilisateur, quel outil appeler et
        avec quels arguments. Cette capacité est exploitée par l'assistant de monitoring
        de notre solution pour interroger l'état de l'infrastructure (métriques
        Prometheus, état des conteneurs Docker, journaux applicatifs) sans que
        l'utilisateur ait à manipuler directement ces interfaces.

% ─────────────────────────────────────────────────────────────────────────────
\section{Orchestration de workflows IA}

    Un pipeline de revue de sécurité combine plusieurs étapes hétérogènes : clonage du
    dépôt, exécution parallèle de scanners, classification par un premier LLM, analyse
    approfondie par un second, publication des résultats. Coordonner ces étapes avec un
    simple script séquentiel devient rapidement intenable lorsque le pipeline doit
    partager un état entre les nœuds, router dynamiquement selon une classification,
    survivre à un redémarrage du conteneur ou s'interrompre pour une approbation humaine.

    L'orchestration de workflows IA répond à ces besoins en modélisant le pipeline sous
    forme de \textbf{graphe d'état dirigé} : chaque nœud représente une étape (souvent
    un appel LLM ou un outil), chaque arête une transition pouvant être conditionnée
    par l'état courant. Trois propriétés distinguent ce paradigme d'un script classique :
    la \textbf{persistance d'état}, qui permet de sauvegarder un \textit{checkpoint}
    après chaque nœud et de reprendre l'exécution exactement là où elle s'était arrêtée ;
    le \textbf{routage conditionnel}, qui permet à un nœud de décider du prochain en
    fonction de son résultat ; et l'\textbf{interruption humaine}, qui suspend le graphe
    en attente d'une validation externe avant de reprendre. Ces trois propriétés sont
    indispensables à notre cas d'usage, où un pipeline de 15 à 25 minutes doit survivre aux
    redémarrages et où certaines décisions critiques peuvent requérir une revue manuelle.

% ─────────────────────────────────────────────────────────────────────────────
\section{Outils d'analyse de sécurité statique (SAST)}

    \subsection{Présentation générale}

        Les outils SAST (Static Application Security Testing) analysent le code source ou
        les artefacts de build sans exécuter le programme. Aucun outil ne couvre à lui seul
        toutes les catégories de risques. Notre solution combine cinq outils complémentaires,
        chacun expert dans sa catégorie, pour assurer une couverture exhaustive. La figure
        \ref{fig:sast_overview} présente ces cinq outils et leurs domaines respectifs.

        \begin{figure}[H]
            \centering
            \frame{\includegraphics[width=0.85\columnwidth]{sast_tools_overview}}
            \caption{Les cinq outils SAST intégrés — couverture complémentaire des risques}
            \label{fig:sast_overview}
        \end{figure}

    \subsection{Les cinq catégories de risques couvertes}

        Notre pipeline couvre cinq catégories de risques de sécurité, chacune nécessitant
        un outil spécialisé. Les \textbf{CVE images et filesystem} concernent les
        vulnérabilités connues affectant les paquets système, les bibliothèques tierces
        et les images de conteneurs utilisés par l'application. Les \textbf{secrets et
        credentials} regroupent les clés API, mots de passe, jetons d'accès et certificats
        accidentellement commités dans le dépôt. Les \textbf{vulnérabilités applicatives
        OWASP} couvrent les patterns dangereux dans le code source : injections SQL,
        désérialisation non sécurisée, chemins traversaux, etc. L'\textbf{Infrastructure-as-Code}
        désigne les erreurs de configuration dans les Dockerfile, fichiers Terraform et
        manifestes Kubernetes. Enfin, les \textbf{dépendances vulnérables} concernent les
        bibliothèques tierces déclarées dans les fichiers de dépendances
        (\texttt{requirements.txt}, \texttt{package.json}, \texttt{go.mod}, etc.) dont une
        version connue présente une faille publiée.

    \subsection{Réduction des tokens SAST}

        Les sorties brutes des scanners sont volumineuses et contiennent de nombreux champs
        redondants. Envoyer le JSON brut au LLM gaspillerait la fenêtre de contexte au
        détriment de l'analyse effective du code. Une étape de nettoyage réduit le volume
        global de 52\% avant transmission au modèle 14B, comme présenté dans le tableau
        \ref{tab:token_reduction}.

        \begin{table}[H]
        \centering
        \begin{tabular}{|m{3.5cm}|m{2.5cm}|m{2.5cm}|m{2cm}|m{3cm}|}
        \hline
        \textbf{Scanner} & \textbf{Avant} & \textbf{Après} & \textbf{Réduction} & \textbf{Technique} \\
        \hline
        Trivy (30 vulnérabilités) & \textasciitilde6 000 chars & \textasciitilde2 500 chars & 58\% & Suppression Description, CVSS, References \\
        \hline
        Gitleaks & \textasciitilde400 chars & \textasciitilde200 chars & 50\% & Champ Match omis (sécurité LLM) \\
        \hline
        Semgrep & \textasciitilde2 000 chars & \textasciitilde1 200 chars & 40\% & Findings INFO agrégés en compteur \\
        \hline
        Checkov & \textasciitilde1 500 chars & \textasciitilde800 chars & 47\% & Champ guideline URL supprimé \\
        \hline
        \textbf{Total SAST} & \textbf{\textasciitilde9 900} & \textbf{\textasciitilde4 700} & \textbf{52\%} & — \\
        \hline
        \end{tabular}
        \caption{Réduction des tokens SAST avant transmission au LLM}
        \label{tab:token_reduction}
        \end{table}

% ─────────────────────────────────────────────────────────────────────────────
\section{Stack applicative et observabilité}

    \subsection{Framework API asynchrone}

        Le cœur de l'agent expose une interface HTTP qui reçoit les webhooks GitHub,
        publie les résultats et sert l'assistant de chat en temps réel. Cette interface
        doit gérer des opérations longues (scans SAST, appels LLM de plusieurs minutes)
        sans bloquer les autres requêtes entrantes, ce qui impose un framework basé sur
        un modèle d'exécution \textbf{asynchrone} via \texttt{asyncio}. Une validation
        automatique des charges utiles entrantes est également nécessaire pour garantir
        l'intégrité des données reçues depuis les webhooks.

    \subsection{Base de données relationnelle}

        Notre architecture confie à une base de données relationnelle deux rôles distincts.
        D'une part, elle sert de \textbf{base de connaissances sécurité} pour stocker
        l'ensemble des revues de \textit{Pull Requests}, des résultats de scanners et des
        profils de risque par dépôt. Cette base alimente ainsi les requêtes de tendances.
        D'autre part, elle assure la \textbf{persistance des checkpoints} de l'orchestrateur
        de workflows. Un pipeline de 15 à 25 minutes peut ainsi survivre à un redémarrage du
        conteneur.

    \subsection{Cache et déduplication}

        Un cache en mémoire est nécessaire pour trois fonctions opérationnelles critiques.
        La \textbf{déduplication des webhooks} évite le double traitement d'un même
        événement lorsque GitHub réémet une notification après timeout. Le
        \textbf{rate limiting} garantit un maximum de pipelines simultanés par dépôt,
        protégeant le moteur d'inférence contre la surcharge. Enfin, le \textbf{cache des
        résultats de scanners} permet de retourner les résultats en quelques millisecondes
        lors d'un nouveau push sur la même branche, sans relancer une analyse complète.

    \subsection{Observabilité}

        Un système exécutant des pipelines de plusieurs minutes, dépendant d'un moteur
        d'inférence local et orchestrant cinq scanners externes, doit être instrumenté
        pour rester opérable. La stack d'observabilité couvre trois besoins distincts :
        la \textbf{collecte de métriques} (durée du pipeline, latence des appels LLM,
        taux d'erreur des scanners, occupation disque), le \textbf{stockage long terme}
        pour permettre l'analyse de tendances sur 90 jours, et la \textbf{visualisation}
        via des dashboards opérationnels. Un mécanisme d'\textbf{alerting} complète
        l'ensemble pour notifier les anomalies (saturation disque, indisponibilité du
        moteur LLM, échecs répétés de scanners) avant qu'elles n'impactent les utilisateurs.

% ─────────────────────────────────────────────────────────────────────────────
\section{Choix technologiques}

    \subsection{Choix du moteur d'inférence LLM}

        Pour un établissement bancaire, le choix entre inférence locale et API cloud est
        fortement contraint par les exigences réglementaires. Le tableau \ref{tab:ollama_vs_cloud}
        présente une comparaison entre les deux approches.

        \begin{table}[H]
        \centering
        \begin{tabular}{|m{3.5cm}|m{5cm}|m{5cm}|}
        \hline
        \textbf{Critère} & \textbf{Ollama (inférence locale)} & \textbf{API cloud (OpenAI, etc.)} \\
        \hline
        Confidentialité & Totale — code reste sur le VPS & Code transmis à des serveurs tiers \\
        \hline
        Conformité bancaire & Compatible avec les exigences de la BCT & Problématique selon la réglementation \\
        \hline
        Coût d'exploitation & Fixe — VPS existant & Variable selon le volume de tokens \\
        \hline
        Disponibilité & 100\% hors ligne & Dépend de la connectivité et du fournisseur \\
        \hline
        Contrôle de version & Modèle figé, reproductible & Version modifiée sans préavis \\
        \hline
        Latence & 3–12 tok/s sur CPU Haswell & Variable selon congestion réseau \\
        \hline
        \end{tabular}
        \caption{Comparaison entre l'inférence locale (Ollama) et les API cloud LLM}
        \label{tab:ollama_vs_cloud}
        \end{table}

        Après l'évaluation comparative présentée dans le tableau \ref{tab:ollama_vs_cloud},
        nous avons retenu Ollama \cite{Ollama} pour l'inférence locale. La confidentialité
        du code source bancaire et la conformité réglementaire vis-à-vis de la Banque
        Centrale de Tunisie imposent que les diffs de code ne quittent jamais
        l'infrastructure interne.

    \subsection{Choix du modèle LLM}

        Afin de choisir le modèle le mieux adapté à notre infrastructure (12 cœurs Intel
        Haswell AVX2, sans GPU), nous avons conduit un benchmark sur une suite de cinq
        requêtes représentatives avec le prompt système complet de 7 577 caractères. Les
        résultats sont présentés dans le tableau \ref{tab:llm_benchmark}.

        \begin{table}[H]
        \centering
        \begin{tabular}{|m{4cm}|m{1.8cm}|m{2cm}|m{2.5cm}|m{2.5cm}|}
        \hline
        \textbf{Modèle} & \textbf{Taille} & \textbf{Vitesse} & \textbf{Précision outils} & \textbf{Décision} \\
        \hline
        \texttt{qwen2.5-coder:7b} & 4,7 Go & \textasciitilde5 tok/s & 80\% & Chat (défaut) \\
        \hline
        \texttt{qwen2.5-coder:14b} & 9,0 Go & \textasciitilde3,2 tok/s & 80\% & Revue PR \\
        \hline
        \texttt{llama3.2:3b} & 2,0 Go & \textasciitilde8 tok/s & 0\% (prompt complet) & Expérimental \\
        \hline
        \texttt{granite3.1-dense:2b} & 1,6 Go & \textasciitilde8,5 tok/s & 0\% & Incompatible \\
        \hline
        \end{tabular}
        \caption{Benchmark des modèles LLM — 5 requêtes, prompt système complet, CPU BTE}
        \label{tab:llm_benchmark}
        \end{table}

        L'analyse des résultats du benchmark révèle des conclusions importantes. En premier
        lieu, \texttt{llama3.2:3b} atteint 0\% de précision avec le prompt système complet,
        car ce dernier consomme environ 2 730 tokens, occupant les deux tiers de la fenêtre
        de contexte de 4 096 tokens et ne laissant qu'un espace résiduel pour le raisonnement.
        En second lieu, \texttt{granite3.1-dense:2b} utilise le schéma propriétaire IBM
        \texttt{\{"tool\_name": ...\}}, incompatible avec notre format d'appel d'outils.
        Enfin, \texttt{qwen2.5-coder:7b} et \texttt{14b} atteignent une précision identique
        de 80\%, mais le modèle 7b est 43\% plus rapide par token. Nous retenons donc une
        architecture à deux modèles : le 7b pour la classification rapide de la
        \textit{Pull Request} (environ 30 secondes) et le 14b pour l'analyse de sécurité
        approfondie (15 à 25 minutes selon la complexité de la PR sur CPU sans GPU).

    \subsection{Choix de l'orchestrateur}

        Le tableau \ref{tab:langgraph_comp} présente une comparaison entre LangGraph
        et les alternatives d'orchestration envisagées.

        \begin{table}[H]
        \centering
        \begin{tabular}{|m{3.5cm}|m{5cm}|m{5cm}|}
        \hline
        \textbf{Critère} & \textbf{LangGraph} & \textbf{Prefect / Celery} \\
        \hline
        Persistance d'état & Native PostgreSQL (\texttt{AsyncPostgresSaver}) & Implémentation custom requise \\
        \hline
        Reprise après crash & Automatique depuis le checkpoint & Partielle selon configuration \\
        \hline
        Interruption humaine & \texttt{interrupt\_before} natif & Complexe à implémenter \\
        \hline
        Routage conditionnel & Arêtes conditionnelles natives & Nécessite des contournements \\
        \hline
        Intégration LLM & LangChain native & Non native \\
        \hline
        \end{tabular}
        \caption{Comparaison entre LangGraph et les alternatives d'orchestration}
        \label{tab:langgraph_comp}
        \end{table}

        Après l'évaluation présentée dans le tableau \ref{tab:langgraph_comp}, nous avons
        choisi LangGraph \cite{LangGraph}, bibliothèque Python construite sur LangChain
        \cite{LangChain}. Sa persistance d'état native sur PostgreSQL garantit qu'un
        pipeline de 15 à 25 minutes survit aux redémarrages du conteneur et reprend exactement
        là où il s'était arrêté. Il s'agit d'une exigence fondamentale pour notre cas
        d'usage.

    \subsection{Choix des outils SAST}

        Le tableau \ref{tab:sast_comp} présente la comparaison effectuée pour le choix de
        chaque outil selon sa catégorie de risque.

        \begin{table}[H]
        \centering
        \begin{tabular}{|m{2.3cm}|m{2cm}|m{4.5cm}|m{3cm}|}
        \hline
        \textbf{Catégorie} & \textbf{Retenu} & \textbf{Justification} & \textbf{Rejeté} \\
        \hline
        CVE images \& filesystem & Trivy & Léger, JSON structuré, base CVE quotidienne & Grype (moins mature) \\
        \hline
        Secrets \& credentials & Gitleaks & 130+ règles, champ \texttt{Match} omettable & TruffleHog (plus lent) \\
        \hline
        SAST OWASP & Semgrep & Règles OWASP épinglées, déterministe & SonarQube (serveur dédié) \\
        \hline
        Infrastructure-as-Code & Checkov & pip-installable, Dockerfile + Terraform & KICS (moins de règles) \\
        \hline
        Dépendances & OSV-Scanner & Base Google OSV, aucune clé API & Snyk (fonctionnalités payantes) \\
        \hline
        \end{tabular}
        \caption{Étude comparative des outils SAST par catégorie de risque}
        \label{tab:sast_comp}
        \end{table}

        Après évaluation des alternatives, la combinaison Trivy \cite{Trivy} + Gitleaks
        \cite{Gitleaks} + Semgrep \cite{Semgrep} + Checkov \cite{Checkov} + OSV-Scanner
        \cite{OSVScanner} offre la couverture la plus exhaustive, avec des formats de
        sortie JSON homogènes facilitant le nettoyage et la réduction des tokens avant
        transmission au LLM.

    \subsection{Choix de la stack applicative}

        Pour le framework API, nous retenons \textbf{FastAPI} \cite{FastAPI}, framework
        Python ASGI haute performance basé sur les annotations de type Python 3.10+.
        Sa gestion native de l'asynchronisme via \texttt{asyncio} est essentielle pour
        exécuter des sous-processus de scanning et des appels LLM en parallèle sans
        bloquer le thread principal. Il assure également une validation automatique des
        payloads via Pydantic. Cette validation préserve l'intégrité des données reçues
        depuis les webhooks GitHub.

        Pour la base de données relationnelle, nous retenons \textbf{PostgreSQL}
        \cite{PostgreSQL}, qui assume les deux rôles décrits précédemment : la base de
        connaissances sécurité et la persistance des checkpoints LangGraph via
        l'\texttt{AsyncPostgresSaver}. Son écosystème mature et son support natif des
        types JSON facilitent le stockage des résultats hétérogènes des scanners.

        Pour le cache et la déduplication, nous retenons \textbf{Redis} \cite{Redis}.
        Son opération atomique \texttt{SET NX} permet une déduplication fiable des
        webhooks, et sa gestion native des TTL convient parfaitement au cache des
        résultats de scanners (TTL d'une heure) et au rate limiting (maximum trois
        pipelines simultanés par dépôt).

    \subsection{Choix de la stack d'observabilité}

        Le tableau \ref{tab:obs_stack} présente les composants de la stack d'observabilité
        et les alternatives considérées.

        \begin{table}[H]
        \centering
        \begin{tabular}{|m{3cm}|m{4cm}|m{5cm}|}
        \hline
        \textbf{Composant} & \textbf{Rôle} & \textbf{Alternative rejetée} \\
        \hline
        Prometheus \cite{Prometheus} & Collecte de métriques, 15 règles d'alerte & DataDog (cloud externe, coût) \\
        \hline
        VictoriaMetrics \cite{VictoriaMetrics} & Stockage long terme 90 jours & InfluxDB (déploiement plus complexe) \\
        \hline
        Grafana \cite{Grafana} & 3 dashboards de visualisation & Kibana (orienté logs, pas métriques) \\
        \hline
        node-exporter \cite{NodeExporter} & Métriques host réelles depuis \texttt{/host/proc} & cAdvisor (namespace Docker uniquement) \\
        \hline
        AlertManager \cite{AlertManager} & Routage alertes → agent → Slack & PagerDuty (payant, cloud) \\
        \hline
        \end{tabular}
        \caption{Stack d'observabilité et alternatives rejetées}
        \label{tab:obs_stack}
        \end{table}

        La combinaison \textbf{Prometheus + VictoriaMetrics + Grafana + node-exporter +
        AlertManager} a été retenue. C'est d'abord la seule pile entièrement open source
        et auto-hébergeable qui couvre les quatre besoins, ce qui prolonge la contrainte
        de confidentialité déjà appliquée au moteur LLM. Les cinq composants ont par
        ailleurs été conçus pour s'intégrer entre eux sans adaptateur. Prometheus émet
        vers AlertManager nativement et Grafana lit indifféremment Prometheus et
        VictoriaMetrics, ce qui nous a évité d'écrire des couches de pontage. Reste le
        cas particulier de \texttt{node-exporter} \cite{NodeExporter} : sur un VPS qui
        utilise le \textit{containerd snapshotter}, c'était la seule option capable
        d'exposer les vraies métriques de l'hôte, cAdvisor \cite{cAdvisor} restant
        aveugle au noyau sous-jacent.

    \subsection{Synthèse des choix technologiques}

        Le tableau \ref{tab:tech_synthesis} récapitule l'ensemble des choix technologiques
        effectués pour notre solution.

        \begin{table}[H]
        \centering
        \begin{tabular}{|m{5cm}|m{4cm}|m{3.5cm}|}
        \hline
        \textbf{Besoin} & \textbf{Technologie retenue} & \textbf{Version} \\
        \hline
        Framework API async & FastAPI & 0.135+ \\
        \hline
        Orchestration workflow IA & LangGraph & 1.1+ \\
        \hline
        Inférence LLM locale & Ollama & latest \\
        \hline
        Modèle de classification & \texttt{qwen2.5-coder:7b} & 4,7 Go \\
        \hline
        Modèle d'analyse sécurité & \texttt{qwen2.5-coder:14b} & 9,0 Go \\
        \hline
        Base de données & PostgreSQL & 16-alpine \\
        \hline
        Cache / Déduplication & Redis & 7-alpine \\
        \hline
        CVE scanning & Trivy & 0.69+ \\
        \hline
        Détection de secrets & Gitleaks & 8.30+ \\
        \hline
        SAST OWASP & Semgrep & 1.157+ \\
        \hline
        IaC scanning & Checkov & 3.2.517+ \\
        \hline
        Dépendances vulnérables & OSV-Scanner & 2.3+ \\
        \hline
        Métriques & Prometheus + VictoriaMetrics & latest \\
        \hline
        Dashboards & Grafana & latest \\
        \hline
        Métriques host & node-exporter & latest \\
        \hline
        Reverse proxy & Nginx & alpine \\
        \hline
        Conteneurisation & Docker + Compose & 29.4.0 \\
        \hline
        \end{tabular}
        \caption{Tableau de synthèse des choix technologiques}
        \label{tab:tech_synthesis}
        \end{table}

% ─────────────────────────────────────────────────────────────────────────────
\section{Analyse et définition des besoins}

    \subsection{Besoins fonctionnels}

        La plateforme couvre cinq fonctions principales, qui suivent le parcours naturel
        d'une \textit{Pull Request} depuis son ouverture jusqu'à son archivage.

        La première est la \textbf{détection automatisée des vulnérabilités} : à l'ouverture
        ou à la mise à jour d'une PR, l'agent clone le dépôt et lance en parallèle les
        scanners retenus pour la classification observée. Sur cette base, il produit une
        \textbf{revue de sécurité combinée} en un unique appel au modèle 14B, structurée en
        quatre éléments : analyse OWASP Top 10, score de risque
        (\textsc{info} → \textsc{critical}), verdict
        (\textsc{approve} / \textsc{request\_changes} / \textsc{block}) et commentaires
        \textit{inline} porteurs de suggestions de correction. Le résultat est ensuite
        \textbf{publié sur GitHub} sous trois formes : un commentaire global dans
        l'onglet \textit{Conversation}, des commentaires \textit{inline} dans l'onglet
        \textit{Files Changed} et un statut de commit positionné pour servir de
        \textit{gate} au CI/CD. Une fois la PR clôturée, la revue, les sorties brutes des scanners et la
        liste des fichiers touchés sont \textbf{persistés dans PostgreSQL}, ce qui alimente
        les analyses de tendance par dépôt et les requêtes ultérieures depuis le chat.
        Cette dernière fonction est précisément le \textbf{chat opérationnel en temps réel} :
        servi en SSE, il pilote dix-neuf outils qui couvrent l'hôte, Docker, Ollama, Redis,
        Prometheus, PostgreSQL et les artefacts SAST.

    \subsection{Besoins non fonctionnels}

        Au-delà des fonctionnalités elles-mêmes, plusieurs contraintes pèsent sur la
        qualité de service attendue de la plateforme.

        La contrainte la plus structurante est la \textbf{confidentialité} : la totalité
        de l'inférence LLM doit s'exécuter en local, sans qu'aucun code source ni diff
        ne quitte l'infrastructure BTE. Il s'agit de la condition de conformité
        réglementaire imposée par le secteur bancaire. Vient ensuite la \textbf{performance} : le
        pipeline complet, de la réception du webhook à la publication de la revue, doit
        tenir en 15 à 25 minutes selon la complexité de la \textit{Pull Request} pour ne pas
        peser sur les cycles de livraison. La
        \textbf{résilience} et la \textbf{fiabilité} sont deux exigences proches mais
        distinctes. La première impose que le système continue à fonctionner si un
        scanner ou le modèle LLM devient temporairement injoignable, ce que prennent en
        charge le circuit breaker et les fallbacks. La seconde impose qu'un pipeline en
        cours d'exécution survive à un redémarrage du conteneur, ce qui repose sur la
        persistance des checkpoints LangGraph dans PostgreSQL. L'\textbf{observabilité}
        demande que toutes les métriques critiques (durée du pipeline, taux d'erreur,
        usage disque, état d'Ollama) remontent à Prometheus et soient consultables dans
        Grafana. Enfin, la \textbf{sécurité de la plateforme elle-même} couvre trois
        points concrets : signature HMAC-SHA256 des webhooks, authentification sur les
        interfaces exposées, et exclusion des secrets détectés par Gitleaks du payload
        envoyé au LLM.

\section*{Conclusion}
\addcontentsline{toc}{section}{Conclusion}

Ce chapitre a posé les bases conceptuelles du projet. Nous avons revisité le DevOps,
la CI/CD et leur extension DevSecOps qui intègre la sécurité à chaque étape du cycle ;
puis le fonctionnement des grands modèles de langage et l'enjeu particulier de
l'inférence locale dans un contexte bancaire ; enfin les principes d'orchestration de
workflows IA et les cinq familles de risques couvertes par les outils SAST. L'étude
comparative qui a suivi a fixé les choix technologiques retenus : Ollama \cite{Ollama}
pour la confidentialité bancaire, LangGraph \cite{LangGraph} pour la persistance native
de l'état, et l'assemblage Trivy / Gitleaks / Semgrep / Checkov / OSV-Scanner pour la
couverture SAST. Les besoins fonctionnels et non fonctionnels recensés en fin de
chapitre cadrent la conception détaillée du chapitre suivant.
